\documentclass{article}
\usepackage{graphicx}
\usepackage{fancyhdr}
\usepackage{geometry}
\usepackage{setspace}
\usepackage{tikz}
\usepackage[italian]{babel}
\usepackage[hidelinks]{hyperref}
\usepackage{float}

% Margini della pagina
\geometry{a4paper, margin=1in}

% Intestazione personalizzata
\pagestyle{fancy}
\fancyhf{}
\fancyhead[L]{Code7Crusaders - Software Development Team}
\fancyhead[R]{\thepage}

% Spaziatura delle righe
\setstretch{1.2}

\begin{document}

% Sezione del titolo
\begin{titlepage}

    \AddToHookNext{shipout/background}{
    \begin{tikzpicture}[remember picture,overlay]
    \node at (current page.center) {
    \includegraphics[width=1.05\paperwidth]{../../img/background.png}
    };
    \end{tikzpicture}
    }

    \centering
    \vspace*{2cm}
    
    \includegraphics[width=0.3\textwidth]{../../img/logo/7Crusaders_logo.png} % Aggiungi il logo qui
    \vspace{1cm}
    
    {\Huge \textbf{Code7Crusaders}}\\
    \vspace{0.5cm}
    {\Large Software Development Team}\\
    \vspace{2cm}
    
    {\large \textbf{Analisi per la scelta delle tecnologie backend}}\\
    \vspace{5cm}

    \textbf{Membri del Team:}\\
    Enrico Cotti Cottini, Gabriele Di Pietro, Tommaso Diviesti \\
    Francesco Lapenna, Matthew Pan, Eddy Pinarello, Filippo Rizzolo \\
    \vspace{0.5cm}
    
    {\large \textbf{Data:}} 12 Febbraio 2025\\
    
    \vspace{1cm}
\end{titlepage}
\clearpage

% Indice
\newpage
\tableofcontents
\newpage

% Sezione Introduzione
\section{Obiettivo}
Questo documento si pone l'obiettivo di confrontare le diverse 
tecnologie backend in modo da avere un'idea precisa e prendere 
una decisione sicura per la scelta di esse e per la loro integrazione 
all'interno del nostro progetto. Il confronto considera aspetti 
tecnici, caratteristiche, vantaggi e svantaggi dei vari framework 
discussi. Le nostre scelte finali sono l'utilizzo di Flask e di 
Langchain per diverse motivazioni chiarite nei paragrafi sucessivi.

\section{Analisi tecnologie backend scelte}

\subsection{Flask}

\subsection{Langchain}




\section{Possibili alternative da valutare}

\subsection{FastAPI}

\subsection{Django REST Framework}

\subsection{LlamaIndex}

\subsection{Haystack}



\section{Confronto}


\section{Conclusioni}
\textbf{Flask} è stato scelto per la sua leggerezza e semplicità nel creare 
API RESTful. Essendo un micro-framework, permette di sviluppare 
rapidamente un backend senza imporre dei vincoli rigidi. 
Inoltre, la sua ampia documentazione e la sua flessibilità lo rendono 
ideale per prototipi e progetti in evoluzione.
Per quanto riguarda \textbf{Langchain}, invece, si tratta di una libreria
progettata per facilitare l'integrazione dei \textit{Large Language Models (LLMs)}
nei sistemi Software. Esso permette di gestire conversazioni, memoria
contestuale e connettori a database vettoriali, rendendo l'interazione 
con i modelli più strutturata e personalizzabile.

\end{document}